\begin{textsample}{POS Dim 3 – 2010\_c6 – Score -6.00 – 2010\_c6/enade\_2010\_med\_vet\_02adfbfccf.txt}  \label{ex:f3_pos_267}
CAPÍTULO 3 ANÁLISE TÉCNICA DA PROVA Este capítulo tem por objetivo apresentar o desempenho dos estudantes de Medicina Veterinária no ENADE/2010.
Para isso, foram calculadas as estatísticas básicas da prova em seu todo, bem como as estatísticas dos componentes relacionadas à Formação Geral e ao de Conhecimento Específico da Área.
Nas tabelas, são apresentados o tamanho da população inscrita e de presentes, e as seguintes estatísticas das notas⁷ : média do desempenho na prova, erro-padrão da média, desvio padrão, nota mínima, mediana, nota máxima, coeficiente de assimetria e de variação.
As estatísticas apresentadas neste capítulo contemplam o total de estudantes e, separadamente, os Ingressantes e os Concluintes.
Foram calculadas tendo -se em vista as seguintes agregações : ( a ) as Grandes Regiões e o país como um todo ; ( b ) a Categoria Administrativa ; e ( c ) a Organização Acadêmica.
Em relação aos gráficos de distribuição de notas, o intervalo considerado foi de 10 unidades, aberto à esquerda e fechado à direita, com exceção do primeiro intervalo, [ 0 ; 10 ], fechado em ambos os extremos.
Para os gráficos de distribuição das notas das questões discursivas, foram consideradas mais duas categorias : questão em branco e nota zero. 3.1 ESTATÍSTICAS BÁSICAS DA PROVA 3.1.1 Estatísticas Básicas Gerais A Tabela 3.1 apresenta as estatísticas básicas da prova por grupo de estudantes.
A população total de estudantes regularmente inscritos⁸ para a prova foi 16.684.
Destes, 14.382 estiveram presentes, sendo 13,8 \% o índice de não comparecimento.
A abstenção foi maior entre Ingressantes ( 19,3 \% ) do que entre Concluintes ( 3,8 \% ).
A nota média de toda a prova foi 39,2, sendo que os Ingressantes obtiveram média mais baixa, 33,2, que os Concluintes, cuja média foi 48,4.
O desvio padrão foi 13,9, sendo o do grupo de Ingressantes menor ( 11,6 ) do que o do grupo dos Concluintes ( 12,0 ) ; os coeficientes de variação ( 34,9 \% no grupo dos Ingressantes e 24,8 \% no dos Concluintes ) indicam menor dispersão relativa nas notas de Concluintes.
A nota máxima foi 81,0, obtida por Concluintes, ao passo que a maior nota obtida por Ingressantes foi 71,6. 7 Algumas dessas estatísticas estão definidas no Capítulo 1. 8 As estatísticas das notas apresentadas consideram apenas os estudantes regularmente inscritos e com resultados validados por ausência de erro administrativo.
Uma comparação entre o desempenho de Ingressantes e Concluintes também pode ser realizada a partir dos dados apresentados no Gráfico 3.1, com a distribuição de notas finais.
Destaca -se que as notas mais elevadas, considerando -se o intervalo posterior à nota 40,0, foram atingidas fundamentalmente pelos Concluintes, cuja concentração das notas está à direita dos Ingressantes.
Já nas faixas de notas mais baixas, inverteu -se o quadro identificado nas faixas de notas mais altas.
Essa característica é corroborada pelo valor da nota mediana 32,9 no grupo dos Ingressantes, abaixo da encontrada no grupo dos Concluintes, 49,4 transparecendo uma superioridade de desempenho dos Concluintes em relação aos Ingressantes na prova em tela ( Tabela 3.1 ).
Os coeficientes de assimetria indicam uma assimetria negativa ( -0,5 ) para os Concluintes, e positiva ( 0,2 ) para Ingressantes.
O Gráfico 3.2 apresenta informações referentes à nota final de Ingressantes e Concluintes, comparando -se os resultados de acordo com as Grandes Regiões do país, a Categoria Administrativa e a Organização Acadêmica.
Considerando -se as notas médias dos alunos Ingressantes em cada Grande Região, observa -se que, em relação às notas aferidas em âmbito nacional ( média de 33,2 para Ingressantes e de 48,4 para Concluintes ) que a menor média foi obtida na região Sudeste ( 32,3 ) e a maior, na região Nordeste ( 34,4 ).
Entre os Concluintes, também existem diferenças estatísticas dignas de comentários.
A maior média de notas foi obtida também pela região Nordeste ( 49,8 ), ao passo que a menor foi identificada na região Norte ( 46,2 ).
As demais médias de Concluintes foram 48,0 para as regiões Sudeste e Centro-Oeste, e 49,3 para a região Sul.
Destaca -se a região Nordeste com médias superiores às nacionais, tanto para o grupo de Ingressantes quanto para o de Concluintes.
Levando -se em conta os grupos dos estudantes em Categorias Administrativas, observa -se que a pontuação média mais elevada entre os Concluintes foi encontrada em instituições públicas ( 52,7 ), com desempenho acima da média nacional.
Já os Concluintes provenientes de instituições privadas obtiveram média 45,8, abaixo dessa média.
Quanto aos Ingressantes, as notas mais elevadas, também acima da média nacional para esse grupo de estudantes, ocorreram nas instituições públicas ( 37,9 ).
Os Ingressantes vinculados a instituições privadas obtiveram nota média 30,1, abaixo da média nacional.
Tendo como foco as notas médias dos estudantes agrupados por Organização Acadêmica, distribuídas entre Universidades, Centros Universitários e Faculdades, para os Concluintes vinculados a Universidades, a média 49,9 ficou acima da média nacional.
Pode -se verificar a mesma realidade entre os Ingressantes de Universidades, cuja nota média foi 35,1, também acima da nacional para esse grupo.
Considerando -se as demais Organizações Acadêmicas, observa -se que nos Centros Universitários e Faculdades, o desempenho de Ingressantes e de Concluintes ficou abaixo da média nacional, com notas médias 46,8 e 43,7, respectivamente.
Ingressantes das mesmas Organizações Acadêmicas obtiveram médias 28,8 e 30,1, respectivamente. 3.1.2 Estatísticas Básicas no Componente de Formação Geral A Tabela 3.2 apresenta as estatísticas básicas em relação ao componente da prova que avalia a Formação Geral dos estudantes.
Os Concluintes obtiveram desempenho médio 47,9, que foi superior ao dos Ingressantes, 41,6.
Quanto à variabilidade das notas nos dois grupos, o desvio padrão das notas obtidas foi 18,3 para Ingressantes e 16,9 para Concluintes.
Os dois grupos alcançaram notas máximas próximas, sendo 94,0 a maior nota obtida por pelo menos um Ingressante, e 97,0 a maior nota entre os Concluintes.
O Gráfico 3.3 propicia a comparação entre Concluintes e Ingressantes relativamente ao desempenho no componente de Formação Geral.
Nesse componente, tal como na prova como um todo, os Concluintes apresentaram desempenho superior.
Dos Ingressantes, 18,2 \% obtiveram notas na faixa entre 40 e 50 pontos, sendo essa a faixa modal.
Nas faixas de notas seguintes, mais elevadas, verifica -se queda gradual de ocorrência de estudantes Ingressantes.
Entre os Concluintes, houve percentual maior de estudantes no intervalo entre 50 e 60 pontos, com 23,4 \%.
Também foi observada queda gradual no desempenho, nos intervalos de notas mais elevados, com 17,3 \% desses estudantes tendo atingido a faixa de 60 a 70 e 1,2 \%, a de notas entre 80 a 90 pontos.
Dos Concluintes, 2,2 \% tiveram desempenho na menor faixa de notas, que vai de zero a 10 pontos.
O coeficiente de assimetria foi negativo : -0,3 para Concluintes e -0,1 para Ingressantes.
No Gráfico 3.4 são apresentadas as informações referentes ao desempenho dos Ingressantes e dos Concluintes no componente de Formação Geral, em diferentes agregações como Grande Região do país, Categoria Administrativa e Organização Acadêmica.
Considerando -se as notas médias dos estudantes segundo cada Grande Região, observa -se que no grupo de Ingressantes existe diferença estatisticamente significativa ao nível de 95 \% entre as médias na região Nordeste ( 45,1, a mais alta ) e Sudeste ( 40,0, a mais baixa ).
No grupo de Concluintes, a nota mais elevada também ocorreu na região Nordeste ( 49,8 ), destacando -se, ainda, a Sul ( 49,2 ).
A diferença entre notas mais altas e mais baixas para Concluintes não foi estatisticamente significativa ao nível de 95 \% : a menor nota foi 46,8, correspondente aos da região Norte.
Destaca -se que a região Sudeste ficou abaixo da média nacional no componente de Formação Geral também para Concluintes ( 47,1 ).
Tendo como foco as Categorias Administrativas, as instituições públicas obtiveram médias mais elevadas do que as privadas : para Concluintes, notas médias de 53,1 em cursos de IES públicas e de 44,7 nos das privadas.
Para os Ingressantes, as notas foram de, respectivamente, 49,9 e 36,4.
Em comparação à média nacional, as notas de Ingressantes e de Concluintes vinculados às instituições públicas foram superiores, ao passo que, nas instituições privadas, foram inferiores.
Considerando -se o tipo de Organização Acadêmica, nas Universidades houve desempenho melhor do que em Centros Universitários e Faculdades.
Neste primeiro tipo de Organização Acadêmica, Ingressantes tiveram nota média 45,1, enquanto Concluintes, 49,7.
Nos Centros Universitários e Faculdades, os valores foram mais baixos : os Ingressantes tiveram média 35,2 nos Centros Universitários e 35,8 nas Faculdades ; entre Concluintes, a média foi 45,1 para Centros Universitários e 42,8 para as Faculdades. 3.1.3 Estatísticas Básicas do Componente de Conhecimento Específico A Tabela 3.3 apresenta as estatísticas básicas referentes à parte do componente de Conhecimento Específico da Área de Medicina Veterinária.
A média do desempenho foi de 37,5.
Os Concluintes apresentaram desempenho médio ( 48,6 ) superior ao dos Ingressantes ( 30,3 ) ; as medidas de dispersão – desvio padrão e amplitude geral das notas – são um pouco menores no grupo dos Ingressantes do que no dos Concluintes.
Os coeficientes de variação, 40,3 \% no grupo dos Ingressantes e 26,3 \% no dos Concluintes, indicam dispersão relativa menor nas notas desses últimos.
As notas máximas obtidas por Ingressantes e Concluintes foram iguais a 73,0 e 83,5, respectivamente, notas inferiores àquelas encontradas no componente de Formação Geral ( 94,0 e 97,0, respectivamente ).
As medianas das distribuições das notas dos Ingressantes e dos Concluintes foram 29,8 e 49,5, respectivamente, enfatizando a superioridade dos Concluintes neste componente da prova.
Assim como os Gráficos 3.1 e 3.3, o Gráfico 3.5, apresentado a seguir, proporciona uma comparação do desempenho entre Concluintes e Ingressantes, em relação ao componente de Conhecimento Específico.
Dentre as três comparações feitas, observa -se que, embora a distribuição das notas seja semelhante à apresentada na prova como um todo, neste componente, a diferença de desempenho entre Concluintes e Ingressantes é mais acentuada, destacando -se o melhor desempenho dos primeiros.
Claramente a distribuição das notas dos Ingressantes, mais concentradas à esquerda em relação aos Concluintes, evidencia o percentual mais alto de notas nas faixas mais baixas para o primeiro grupo em comparação com o segundo.
A maior concentração de Concluintes está no intervalo de ( 50 a 60 ] pontos com 30,3 \% dos estudantes do grupo, ao passo que a de Ingressantes está na faixa ( 20 a 30 ] pontos com 35,5 \%.
Assim como os Gráficos 3.2 e 3.4, o Gráfico 3.6 apresenta uma comparação dos resultados em relação à Grande Região do país, à Categoria Administrativa e à Organização Acadêmica, agora levando em conta o desempenho dos Ingressantes e dos Concluintes no componente de Conhecimento Específico da prova.
Ao se observar as notas médias dos estudantes, segundo cada Grande Região do país, verifica -se que, entre os Concluintes, a mais elevada foi encontrada na região Nordeste ( 49,8 ), maior que a média nacional ( 48,6 ).
Em relação aos Ingressantes, a região Sul apresenta a média mais elevada ( 31,1 ), seguida da região Nordeste, com média igual a 30,8.
As médias mais baixas entre os Concluintes ( 46,0 ) e entre os Ingressantes ( 29,4 ) foram verificadas na região Norte.
Quanto à Categoria Administrativa, observa -se comportamento semelhante àquele da parte relativa à Formação Geral e à prova como um todo, ou seja, as médias mais elevadas tanto entre os Concluintes quanto entre os Ingressantes foram encontradas em instituições de ensino públicas.
Nesse tipo de IES, a nota média de Concluintes foi 52,5, enquanto a de Ingressantes, 34,0.
Tais médias estão situadas acima da média nacional.
Nas instituições privadas, as médias foram, mais uma vez, inferiores em relação às públicas.
Concluintes obtiveram média 46,1 e Ingressantes, 28,0.
Em relação à Organização Acadêmica, o desempenho dos Concluintes foi superior ao dos Ingressantes nas três categorias.
Nas Universidades, os Concluintes obtiveram a média mais elevada, igual a 50,0.
Um pouco abaixo estiveram os Centros Universitários ( 47,3 ) seguidos pelas Faculdades ( 44,0 ).
Os Ingressantes de Universidades também tiveram nota um pouco mais elevada ( 31,8 ).
Os Ingressantes de Faculdades ( 28,2 ) obtiveram média mais elevada do que os de Centros Universitários ( 26,7 ). 3.2 ANÁLISE DAS QUESTÕES OBJETIVAS 3.2.1 Componente de Formação Geral A Tabela 3.4 apresenta as estatísticas básicas relativas às oito questões objetivas do componente da prova que abrange a Formação Geral dos estudantes.
Como pode ser observado, os estudantes Concluintes obtiveram desempenho superior, com nota média 51,9, em relação aos alunos Ingressantes, cuja nota média foi 45,4.
As medidas de dispersão nos dois grupos de alunos foram similares, como indicam os desvios padrões, assim como a amplitude total das notas, sendo que a nota mínima foi zero e a máxima, 100,0, para os dois grupos de alunos avaliados.
Os coeficientes de variação, 42,7 \% no grupo dos Ingressantes e 34,7 \% no dos Concluintes, indicam menor dispersão relativa nas notas desses últimos.
A mediana foi igual a 50,0 para os dois grupos.
O coeficiente de assimetria do total de estudantes revelou assimetria negativa ( -0,2 ).
A Tabela 3.5 apresenta as classificações das questões objetivas do componente de Formação Geral, segundo o índice de facilidade.
As questões objetivas da prova foram assim avaliadas : das oito questões, nenhuma teve o índice de facilidade classificado como muito fácil.
Quatro questões foram tidas como fáceis, por terem índice de acertos situado na faixa entre 0,61 e 0,85 ( de 61,0 \% a 85,0 \% de acertos ).
Uma questão foi considerada média, situando -se no intervalo entre 0,41 e 0,60 do índice de facilidade, ou seja, houve entre 41,0 \% e 60,0 \% de acertos.
Também uma questão foi classificada como difícil, situando -se no intervalo entre 0,16 e 0,40.
Por fim, duas questões apresentaram menos de 0,15 de acertos, razão pela qual foram classificadas como muito difíceis.
A Tabela 3.6 descreve os resultados das análises das questões objetivas relativas à Formação Geral, segundo o poder de discriminação, utilizando -se para tal, do índice de discriminação ponto bisserial.
Cinco das oito questões apresentaram índices acima de 0,40 e, assim, foram classificadas como muito boas nesta característica para esse grupo de alunos ; uma questão teve um bom índice de discriminação, entre 0,30 e 0,39.
Foram classificadas como de média discriminação duas das questões.
Nenhuma teve fraco nível de discriminação.
Na sequência, na Tabela 3.7, estão especificadas, por questão, os resultados da análise relativa aos índices de facilidade e de discriminação de cada questão no componente relativo à Formação Geral da prova de Medicina Veterinária.
O índice de facilidade variou de 0,12 a 0,80, e o de discriminação, de 0,25 a 0,55.
As questões com índices de discriminação muito bom, de números 1, 3, 4, 5 e 8, figuraram entre as mais fáceis desse conjunto, sendo que quase todas foram classificadas na classe fácil segundo o índice de facilidade.
A exceção foi a questão 1 que foi classificada como média.
As questões 5 e 8 foram as que apresentaram maior poder discriminatório, ambas com índice 0,55, e foram também as mais fáceis, com uma proporção, no caso da questão 8, de 0,80 de acertos.
As questões de números 6 e 7 apresentaram índice de discriminação médio, com valores, respectivamente, de 0,29 e 0,25.
Foram também as mais difíceis, com índice de facilidade, respectivamente, 0,12 e A Tabela 3.6 descreve os resultados das análises das questões objetivas relativas à Formação Geral, segundo o poder de discriminação, utilizando -se para tal, do índice de discriminação ponto bisserial.
Cinco das oito questões apresentaram índices acima de 0,40 e, assim, foram classificadas como muito boas nesta característica para esse grupo de alunos ; uma questão teve um bom índice de discriminação, entre 0,30 e 0,39.
Foram classificadas como de média discriminação duas das questões.
Nenhuma teve fraco nível de discriminação.
Na sequência, na Tabela 3.7, estão especificadas, por questão, os resultados da análise relativa aos índices de facilidade e de discriminação de cada questão no componente relativo à Formação Geral da prova de Medicina Veterinária.
O índice de facilidade variou de 0,12 a 0,80, e o de discriminação, de 0,25 a 0,55.
As questões com índices de discriminação muito bom, de números 1, 3, 4, 5 e 8, figuraram entre as mais fáceis desse conjunto, sendo que quase todas foram classificadas na classe fácil segundo o índice de facilidade.
A exceção foi a questão 1 que foi classificada como média.
As questões 5 e 8 foram as que apresentaram maior poder discriminatório, ambas com índice 0,55, e foram também as mais fáceis, com uma proporção, no caso da questão 8, de 0,80 de acertos.
As questões de números 6 e 7 apresentaram índice de discriminação médio, com valores, respectivamente, de 0,29 e 0,25.
Foram também as mais difíceis, com índice de facilidade, respectivamente, 0,12 e 0,13.
Em comparação às demais, um quantitativo menor de estudantes conseguiu resolvê -las, dentro do universo de participantes.
O Gráfico 3.7 analisa o comportamento da questão de número 8 de Formação Geral.
Trata -se da questão mais fácil e com o maior índice de discriminação desse componente da prova.
No gráfico a seguir, cada uma das cinco curvas representa o percentual de respostas em determinada alternativa da questão, em função da nota dos estudantes nesta parte da prova ( Formação Geral/Múltipla Escolha ) antes de possíveis eliminações pelo critério do ponto bisserial.
A curva em rosa corresponde à alternativa B, a correta para este item.
Assim, observa -se que entre os estudantes com notas mais baixas, ou seja, com desempenho mais fraco neste componente do exame, a situação mais frequente foi a escolha de uma das alternativas incorretas : a alternativa C ( em verde ) ou E ( em roxo ).
Na medida em que a nota aumenta, indicando desempenho melhor nesta parte da prova, aumenta, concomitantemente, a proporção de estudantes que selecionaram a alternativa correta B, atingindo 100 \% para as notas acima de 80,0.
Essa análise permite verificar como a questão discriminou entre os grupos de desempenho, justificando o alto índice obtido na questão.
Os gráficos relativos às demais questões de Formação Geral constam do Anexo I. 3.2.2 Componente de Conhecimento Específico A Tabela 3.8 apresenta as estatísticas básicas em relação às questões objetivas do componente de Conhecimento Específico da prova.
Nesse conjunto de questões, os estudantes Concluintes ( 52,6 ) obtiveram nota média superior, aos Ingressantes ( 34,9 ).
Quanto à variabilidade das notas, o desvio padrão no grupo dos Ingressantes foi 13,9, e no grupo dos Concluintes, 13,5.
Os coeficientes de variação foram 39,8 \% e 25,7 \%, respectivamente, indicando menor dispersão relativa nas notas dos Concluintes.
Com relação à mediana, no grupo dos Ingressantes, esta foi igual a 35,0, da superando a média desse grupo, enquanto no grupo dos Concluintes, a mediana foi igual a 55,0, também um pouco acima da média desse grupo.
O coeficiente de assimetria foi 0,2 para Ingressantes e negativo ( -0,4 ) para Concluintes. 3.2.2 Componente de Conhecimento Específico A Tabela 3.8 apresenta as estatísticas básicas em relação às questões objetivas do componente de Conhecimento Específico da prova.
Nesse conjunto de questões, os estudantes Concluintes ( 52,6 ) obtiveram nota média superior, aos Ingressantes ( 34,9 ).
Quanto à variabilidade das notas, o desvio padrão no grupo dos Ingressantes foi 13,9, e no grupo dos Concluintes, 13,5.
Os coeficientes de variação foram 39,8 \% e 25,7 \%, respectivamente, indicando menor dispersão relativa nas notas dos Concluintes.
Com relação à mediana, no grupo dos Ingressantes, esta foi igual a 35,0, da superando a média desse grupo, enquanto no grupo dos Concluintes, a mediana foi igual a 55,0, também um pouco acima da média desse grupo.
O coeficiente de assimetria foi 0,2 para Ingressantes e negativo ( -0,4 ) para Concluintes.
A Tabela 3.9 apresenta as classificações das questões objetivas da parte da prova relativa ao componente de Conhecimento Específico, segundo o índice de facilidade.
Deve -se considerar o fato de as questões de números 17 e 33 terem sido anuladas pela Comissão, de modo que a classificação foi estabelecida com base em 25 das 27 questões.
A partir dos índices obtidos, pode -se concluir que, quase metade das questões objetivas da prova foram tidas como difíceis : das 25 questões, nove foram classificadas como difíceis e três como muito difíceis.
Não houve questões classificadas como muito fáceis, ao passo que, duas foram tidas como fáceis, na faixa de 0,61 a 0,85 do índice de facilidade, e onze consideradas médias, entre 0,41 e 0,60.
Na Tabela 3.10 são apresentados os índices de discriminação das questões objetivas do componente de Conhecimento Específico da prova.
Observa -se, segundo os índices estabelecidos, que a capacidade de discriminação das questões objetivas desse componente da prova foi classificada, para nove das 25 questões válidas, como boa, enquanto cinco delas tiveram índice de discriminação muito bom.
A maioria das questões – 14 em 25 – proporcionou, portanto, índices de discriminação bons ou muito bons.
Seis delas foram classificadas como médias e cinco como fracas, sendo onze, por conseguinte, a quantidade de questões nos dois patamares mais baixos de discriminação.
Constata -se, assim, que a prova – no que se refere ao componente de Conhecimento Específico – possuía capacidade relativamente boa de discriminar, entre aqueles que dominam ou não o conteúdo.
As questões 12, 13, 22, 24 e 25 foram eliminadas pelo critério do índice de discriminação e não foram utilizadas no cômputo da nota final.
Na Tabela 3.11 estão especificados, por questão, os resultados da análise relativa aos índices de facilidade e de discriminação de cada uma das questões do componente de Conhecimento Específico da prova de Medicina Veterinária.
Dentre as questões que alcançaram os maiores índices de discriminação, muito bom, as questões de números 31 e 37 foram classificadas com respeito ao índice de facilidade como fáceis e as questões 16, 19 e 34 como de dificuldade média.
Em termos de proporção, tais questões variaram de 0,49 a 0,66 quanto aos acertos.
Seus índices de discriminação situaram -se de 0,40 a 0,48.
A questão de número 13 foi a mais difícil dentre as 25 questões específicas válidas, com baixo índice de facilidade, apenas 0,11, o que, em termos percentuais, corresponde a apenas 11,0 \% de acertos.
Essa questão apresentou poder discriminatório igualmente baixo, 0,05, o que comprova ter sido esta a mais difícil para os estudantes.
Como já comentado, foi eliminada do cálculo da nota final.
Destacam -se, também, as questões de números 12 e 22 para as quais os índices de facilidade foram 0,14 e 0,13, respectivamente, o que, em termos percentuais, corresponde a 14,0 \% e 13,0 \% de estudantes que responderam acertadamente.
Tais questões foram, portanto, duas das mais difíceis da prova.
Obtiveram índice de discriminação fraco apenas 0,02 para a questão 12 e 0,10 para a questão 22 e foram também eliminadas.
O Gráfico 3.8 analisa o comportamento da questão 36 do componente de Conhecimento Específico.
Essa questão apresentou índice de facilidade 0,49, ou seja, 49,0 \% dos estudantes assinalaram acertadamente a opção B, correspondente ao gabarito.
Seu índice de discriminação foi igual a 0,39, classificado como bom, embora outras questões que apresentaram percentuais mais elevados de respostas corretas tenham tido índices de discriminação também mais elevados, como o caso da questão de número 34, 0,59 de índice de facilidade e 0,48 de índice de discriminação.
Neste gráfico, cada uma das cinco curvas representa o percentual de respostas em determinada alternativa da questão 36, em função da nota dos estudantes nesta parte da prova antes de possíveis eliminações de questões pelo critério do ponto bisserial.
A alternativa correta B, representada no gráfico pela curva em rosa, foi escolhida em maiores proporções pelos alunos com desempenho melhor nesse componente do exame, enquanto as alternativas incorretas, também denominadas distratores, foram selecionadas principalmente por aqueles com notas mais baixas, ainda que aqueles com nota zero, na sua quase totalidade tenham deixado em branco ou marcaram mais de uma alternativa, comportamento considerado inválido.
A proporção de alunos que selecionaram a resposta correta B aumenta gradativamente com a nota dos alunos chegando a atingir 100 \% para notas acima de 80,0, enquanto a proporção dos que escolhem alternativas incorretas, decai a partir da primeira nota não nula como função da nota nesta parte da prova.
Os gráficos relativos às demais questões do componente de Conhecimento Específico constam do Anexo I. 3.3 ANÁLISE DAS QUESTÕES DISCURSIVAS 3.3.1 Componente de Formação Geral A análise dos resultados de desempenho dos estudantes de Medicina Veterinária, nas duas questões discursivas relativas à Formação Geral, encontra -se na Tabela 3.12 e no Gráfico 3.9.
Observa -se, a partir destes, que as notas médias foram mais baixas nesse conjunto que no conjunto das objetivas.
A média geral dos Ingressantes nas questões objetivas de Formação Geral foi 45,4 ( Tabela 3.4 ), ao passo que nas questões discursivas, a média foi 35,9 para esse grupo.
O mesmo ocorreu entre os Concluintes, que tiveram, em Formação Geral, média 51,9 nas questões objetivas e 42,0 nas questões discursivas.
Cabe destacar que, nas questões discursivas, o desempenho dos Concluintes foi superior ao dos Ingressantes, embora os dois grupos tenham atingido a mesma nota máxima, 100,0.
Os coeficientes de assimetria indicam assimetria negativa, de -0,2 para os Concluintes, e positiva, de 0,1 para os Ingressantes.
Pode -se observar as distribuições das notas desses grupos, no Gráfico 3.9, apresentado a seguir.
Na sequência, os resultados verificados para cada uma das questões discursivas de Formação Geral serão apresentados, estabelecendo -se relações com os conteúdos englobados em cada uma.
Os comentários da banca corretora a respeito do observado nas respostas dos estudantes, suas impressões e conclusões serão expostos junto à análise de cada questão.
Cumpre esclarecer que, tendo em vista que as questões discursivas de Formação Geral são padronizadas, ou seja, constam de todas as provas, os comentários da banca são os mesmos para todas as carreiras acadêmicas, sendo direcionados a todos os estudantes que participaram do ENADE/2010.
A seguir, serão analisados os desempenhos de Ingressantes e Concluintes da Área de Medicina Veterinária nas duas questões discursivas de Formação Geral do ENADE/2010. 3.3.1.1 Análise da Questão Discursiva 9 do Componente de Formação Geral Os dados obtidos a partir das respostas à questão 9 encontram -se na Tabela 3.13 e no Gráfico 3.10.
Nessa questão – de melhor desempenho dentre as duas de Formação Geral – os Ingressantes tiveram média 39,0, e os Concluintes, 44,9, resultando com diferença de 5,9 pontos entre os dois grupos.
Houve semelhança na distribuição de notas dos dois grupos, observando -se, no entanto, percentual mais elevado de Ingressantes que deixaram a questão em branco, correspondente a 26,9 \%, se comparado aos 19,2 \% de Concluintes.
A semelhança também foi verificada entre aqueles que obtiveram nota zero na questão ainda que não tivessem deixado a questão em branco, cujo percentual foi 1,7 \% em ambos os grupos.
A partir de 40 pontos é possível notar a maior incidência de Concluintes nas faixas denotas.
Entre Ingressantes, a nota mediana foi 45,0 pontos, enquanto Concluintes obtiveram 50,0.
Para os dois grupos a maior proporção de notas nas faixas não nulas ocorre no intervalo de 60 a 70, com 13,1 \% dos Ingressantes e 14,9 \% dos Concluintes.
A nota máxima, 100,0 pontos, foi alcançada tanto pelos Concluintes quanto pelos Ingressantes. 3.3.1.2 Comentários sobre as respostas à Questão Discursiva 9 No que tange mais precisamente à atualidade do tema proposto e ao desenvolvimento do conteúdo, deve ser ressaltado que se mostrou sintonizado com a realidade política brasileira, bem como com a necessária reflexão sobre o papel na sociedade do profissional formado pelas Instituições de Ensino Superior do país.
Dissertar sobre democracia e ética foi algo apropriado para a época em que a prova foi aplicada, pois coincidiu com o período das eleições, em que tais temas são mais discutidos na vida cotidiana.
Não se pode, entretanto, deixar de se levar em conta a tendência de os estudantes terem limitado a associação do tema “ ética ” ao assunto “ política ”, considerando -se que a ética abarca não apenas a política.
Tal associação atingiu, em alguma medida, o alcance mais preciso das respostas.
As respostas apresentadas refletiram motivação para escrever sobre o tema da questão.
Alguns estudantes chegaram a desenvolver textos eloquentes e engajados.
Em alguns casos, o conceito de democracia, por exemplo, foi bem desenvolvido ao longo do texto, simultaneamente às exemplificações pedidas.
Além disso, muitos estudantes associaram o exercício da ética ao seu futuro ofício, reconhecendo o seu papel como cidadão no desempenho de sua profissão.
Essa associação insinuou, também, o esforço de reconhecimento da noção de ética prática, ainda que de maneira codificada.
Um dos pontos positivos revelou -se no fato de, na maioria das respostas analisadas, haver uma preocupação em definir a sociedade democrática.
Isso facilitou a correção e elevou a média de notas.
Outro aspecto positivo foi a presença pertinente de exposições coerentes de comportamento ético de um futuro profissional.
A articulação dos itens “ a ” e “ b ” da questão produziu, portanto, todo tipo de reflexão acerca de uma postura ética na sociedade, inclusive, corporativa. dos itens “ a ” e “ b ” da questão produziu, portanto, todo tipo de reflexão acerca de uma postura ética na sociedade, inclusive, corporativa.
A comprovação de que o tema ética tem sido abordado nos cursos superiores foi outro destaque positivo detectado durante as correções.
Assim, além do caráter crítico de que muitas das respostas corretas se revestiram, os relatos constantes sobre desdobramentos da ética comprovaram a abordagem sobre o tema, no mínimo, a ética profissional, nas carreiras universitárias que participaram do ENADE.
A despeito de alguns estudantes se posicionarem contra o exame, seja anulando a questão ou escrevendo algum comentário desabonador sobre o processo, foi observado que significativa parte dos participantes tem consciência político-acadêmica e percebe por que está realizando aquele tipo de prova.
Desse modo, a maioria dos estudantes tinha algo a “ dizer “ sobre o tema da questão, isto é, valia -se da oportunidade para se posicionar a respeito do assunto trazido à baila.
Ainda do ponto de vista do conteúdo desenvolvido nas respostas, observou -se preocupação solidária e consciência cidadã dos estudantes.
Mais um destaque positivo crucial se refere à habilidade textual.
Os estudantes, por exemplo, expressaram atitudes éticas e não éticas e as contrapunham ao conceito de democracia.
Há casos em que os alunos definiam a corrupção política como um exemplo de atitude não ética, enfatizando a real missão desse cargo em uma sociedade democrática : “ atender aos anseios da população e zelar pelo bem comum ”, escreveram vários estudantes.
No tocante à produção de texto dissertativo, à objetividade e ao caráter sintético, muitos estudantes corresponderam ao solicitado, elaborando pequenos textos dissertativos.
Nessas respostas, percebeu -se demonstração de domínio do assunto focado e utilização dos textos de referência como base para a elaboração.
O enunciado da questão 9 determinou que fosse elaborado um texto dissertativo que contivesse a )..., b )... e c )..., fornecendo um encaminhamento para as respostas que conduziu tanto a um resultado positivo quanto a um negativo, a saber : foi produtiva a resposta dos estudantes no sentido de que conseguiram, na maior parte das vezes, atender plenamente aos três quesitos apresentados na questão ; por outro lado, mais especificamente devido ao uso das três letras indicativas ( a, b e c ), houve tendência à elaboração de respostas divididas em itens, em vez de um texto eminentemente dissertativo.
Muitos estudantes apresentaram respostas nesse formato fragmentado para satisfazer a formalidade do roteiro sugerido no comando da questão.
A ideia-chave que deveria permear toda a resposta – democracia – se perdia em fragmentos desconexos entre si.
Se, por um lado, o roteiro do comando da questão facilitou o raciocínio de alguns dos estudantes, por outro, induziu grande parte deles a dois caminhos : i ) o das respostas curtas e objetivas, do tipo “ um profissional de saúde não ético é aquele que conta para as outras pessoas o que foi feito na consulta com o seu paciente ” e ii ) o outro caminho, que é do contraponto entre o conceito de sociedade democrática e a realidade brasileira.
Esse segundo caminho levou vários universitários a produzirem textos bem escritos, distantes, entretanto, do tema exigido. curtas e objetivas, do tipo “ um profissional de saúde não ético é aquele que conta para as outras pessoas o que foi feito na consulta com o seu paciente ” e ii ) o outro caminho, que é do contraponto entre o conceito de sociedade democrática e a realidade brasileira.
Esse segundo caminho levou vários universitários a produzirem textos bem escritos, distantes, entretanto, do tema exigido.
No que se reporta às especificidades do desenvolvimento do conteúdo, a dificuldade ou falta de competência comunicativa revelada em grande parte das respostas denuncia a pouca exposição dos alunos à leitura de textos escritos de mais excelência.
Tal dificuldade mostrou -se mais evidente na impropriedade no uso de conectores e na incapacidade de elaboração de texto dissertativo, no tratamento de questões mais abstratas.
Muitos estudantes demonstraram não saber em que consistia a ética e confundiram exemplificação com definição : ao tentarem dar exemplos, via de regra, elaboravam uma definição incipiente do que se pedia.
Portanto, a dificuldade dos estudantes em argumentar e articular as ideias com coesão e coerência ainda são os grandes “ vilões ” desse processo avaliativo.
O excesso de fragmentos e textos prévios acerca do tema permitiu que a elaboração das respostas se resumisse a uma colagem ou paráfrase mal feita dos materiais disponíveis.
Em vez de subsidiar as respostas, esse material acabou por influenciar negativamente o desempenho dos estudantes.
Desse modo, verificou -se um nível de letramento ainda significativamente baixo, visto que a maioria apresentou uma visão bastante reducionista e simplificada sobre o tema em pauta.
De um modo geral, ao conceituarem ética e sociedade democrática, os estudantes demonstraram desconhecer o limite entre os dois conceitos, sem aludirem à interseção de ambos.
Foi expressiva a quantidade de redações nas quais questionavam a democracia pelo simples fato de estarem fazendo, a fortiori, aquela prova especificamente.
Em muitas das respostas, os estudantes se diziam incomodados pelo fato de estarem obrigados a realizar o exame, demonstrando sua insatisfação.
Por diversas vezes, os estudantes não foram capazes de sintetizar e de selecionar seu discurso o suficiente para atender ao que era pedido, o que ocasionou a falta de, pelo menos, um dos aspectos da resposta.
Foi constatada grande parte de respostas com problemas graves em termos textuais e gramaticais, em relação ao que se esperaria de pessoas que estão vinculadas ao Ensino ( Técnico ) Superior.
Alguns estudantes sequer conseguiam elaborar uma frase coerente e coesa, apresentando respostas que poderiam ser classificadas como de analfabetismo funcional.
De maneira ampla, manifestou -se da parte dos estudantes uma grande dificuldade no trato da língua escrita, dificuldade essa que se traduz, sobretudo, por um desrespeito bastante frequente às regras relativas : i ) à ortografia ( por exemplo : “ proficionais ”, “ infligir as leis ”, “ corrupitos ”, “ democrácia ”, “ opitar ”, “ manter o cizilo proficionau ” ), troca de “ sociedade democrática ” por “ sociedade demográfica ” e de “ ética ” por “ étnica ” ; ii ) à pontuação ; iii ) à concordância verbal ( “ os político deve (... ) eles pode ” ) ; iv ) ao modo de organização textual ( sem distinção entre o que seja dissertativo, argumentativo ou narrativo ). as leis ”, “ corrupitos ”, “ democrácia ”, “ opitar ”, “ manter o cizilo proficionau ” ), troca de “ sociedade democrática ” por “ sociedade demográfica ” e de “ ética ” por “ étnica ” ; ii ) à pontuação ; iii ) à concordância verbal ( “ os político deve (... ) eles pode ” ) ; iv ) ao modo de organização textual ( sem distinção entre o que seja dissertativo, argumentativo ou narrativo ).
Frisa -se que esse comentário não se restringe ao que se costuma chamar de “ língua padrão ”, mas ao que poderia ser chamado de “ excelência textual ”, que não se mede exclusiva ou enfaticamente pela correção gramatical, mas pela competência, objetividade, clareza e, em alguma medida, pela criatividade no ato de redigir.
Por fim, destaca -se um dos problemas dignos de nota na formulação da questão.
A transcrição do Dicionário Houaiss apresentava defeitos.
Por exemplo, a definição 5 de “ democracia ” indicava que esse é o tipo de “ governo que acata a vontade da maioria, embora respeitando os direitos e a livre expressão das minorias ”.
O emprego de um valor concessivo na oração adverbial, por oposição a um valor aditivo, contradiz as concepções efetivamente democráticas apresentadas de 1 a 4.
Repercussões desse valor concessivo, mal empregado, foram encontradas em incontáveis respostas, as quais diziam ser típico de um governo democrático, seguir a vontade da maioria, “ apesar ” de respeitar as minorias.
Outro desses problemas dignos de nota refere -se à indicação do dicionário quanto a “ governo ”, ao passo que o comando focava a ideia mais ampla de “ sociedade democrática ”. 3.3.1.3 Análise da Questão Discursiva 10 do Componente de Formação Geral O desempenho dos estudantes na questão 10 foi inferior ao obtido na de número 9, com nota média dos Ingressantes de 32,8, e de Concluintes, de 39,1.
A nota máxima também foi atingida em ambos os grupos, sendo que o percentual de alunos que deixaram a questão em branco foi novamente mais elevado para Ingressantes, com 32,8 \%, sendo 24,0 \% a proporção de Concluintes que o fizeram.
A questão discursiva 10, por conseguinte, teve maiores percentuais de estudantes que não a responderam, comparando -se à questão de número 9.
A mediana foi 32,3 para Ingressantes e 32,2 para Concluintes, notas inferiores às da questão anterior.
Houve maior percentual de alunos, considerando -se a distribuição de notas não nulas, no intervalo entre 40 e 50 pontos, tanto para Ingressantes quanto para Concluintes, de 14,0 \% para Ingressantes e 16,6 \% para Concluintes.
Não se manifestou uma superioridade expressiva de Concluintes nas diferentes faixas de notas posteriores a esse intervalo, com equilíbrio no desempenho dos dois grupos, só ligeiramente superior no intervalo de notas entre 70 e 80.
Essas informações estão ilustradas na Tabela 3.14 e no Gráfico 3.11. 3.3.1.4 Comentários sobre as respostas à Questão Discursiva 10 Inicialmente, será apresentada uma panorâmica sobre as respostas, com destaque para os aspectos positivos observados.
Em relação ao desenvolvimento do tema e à vinculação do mesmo à vida cotidiana dos estudantes que participaram do exame, a visão positiva em relação à Educação de Jovens e Adultos ( EJA ) reforça a ideia de que a educação pode trazer crescimento não somente profissional, mas, sobretudo, para ampliar os horizontes da cidadania.
Foram frequentes as respostas que apontavam a educação como a melhor saída, senão a única, para o resgate da cidadania.
A atuação dos governos, em especial o governo federal, foi um ponto destacado nas respostas, as quais sinalizavam para a necessidade de as instâncias governamentais assumirem um papel cada vez mais ativo na consolidação da educação – pública, sobretudo – no país.
As políticas públicas focadas na educação foram, então, um item bastante frequente nas respostas, seja para enaltecer, criticar ou requisitar a ação governamental.
Outro ponto positivo destacado diz respeito à consciência ou ao reconhecimento da precariedade de condições de trabalho dos professores, em geral, no Brasil e, em particular, daqueles ocupados com turmas de jovens e adultos.
Nesse sentido, foram inúmeras as propostas apresentadas na segunda parte das respostas com relação à imperiosidade de melhoria salarial para os docentes, e à melhoria das condições gerais de trabalho desses profissionais.
A associação direta entre valorização do professor e sucesso escolar dos estudantes foi um ponto complementar a essa preocupação com os docentes.
Os estudantes demonstraram uma visão de vínculo no processo de ensino-aprendizagem, no qual alunos e professores são partes inalienáveis entre si.
A questão da repetência e da evasão escolar, referenciadas no suporte da questão, foram então, abordadas pelos alunos por meio do viés dessa valorização docente.
Ainda que os temas específicos da repetência e da evasão não fossem contemplados a contento na maior parte das respostas, ao menos, nessa associação com a qualificação e a valorização dos professores, algo foi trabalhado analiticamente pelos estudantes.
Num plano mais abrangente, houve reconhecimento de que, no processo ensino-aprendizagem, não há atores isolados : alunos, profissionais de educação, governos e a própria sociedade como corresponsáveis pelo sucesso ou fracasso de tal processo.
A atualidade do tema e a sua explícita vinculação com o cotidiano de muitos estudantes que realizavam o exame fizeram com que os respondentes demonstrassem um visível grau de envolvimento com o tema proposto.
Ou seja, foi possível detectar um empenho notório dos estudantes ao redigir suas respostas com seriedade e comprometimento relativo à delicada problemática dos milhões de estudantes trabalhadores do país.
Utilizar todo o espaço destinado à resposta, procurando aproveitar ao máximo a oportunidade de emitir um argumento ou uma opinião revelou -se como uma das evidências do aludido empenho.
A EJA também foi reconhecida como um dos caminhos concretos para o acesso à Universidade, no médio prazo.
As respostas se voltaram, muitas vezes, à condição de depoimentos pessoais sobre a conciliação de trabalho e estudo.
Nessa direção, ao assumirem o papel de depoentes, os estudantes expuseram suas histórias de vida, suas expectativas e percalços, suas esperanças e entraves cotidianos, seus anseios mais ambiciosos e problemas estruturais, que poderiam se resolver por meio da educação.
Nesse inciso, ficava clara a relação que se formava entre educação e mobilidade social, mediatizada por iniciativas governamentais como a EJA.
Uma sugestão plausível e válida para a solução dos problemas com a EJA no país reportou -se à necessidade de maior participação do público alvo na definição de estratégias para essa modalidade de ensino.
Nesse detalhe, era notória a associação das preocupações relativas à questão 9, que trata do problema da democracia e da cidadania, com o problema específico da educação escolar.
Houve um número significativo de sugestões ou propostas para se combater o problema da repetência e da evasão, destacando -se melhores salários para os professores e aulas mais dinâmicas.
A lista dessas sugestões, contudo, era bem mais ampla e diversificada, incluindo ainda : a redução da jornada de trabalho dos jovens e adultos que participam da EJA ; melhoria de infraestrutura nas instituições responsáveis pela EJA ; incentivo financeiro aos estudantes trabalhadores por meio de um sistema de bolsas de estudo ; realização de campanhas motivacionais para o retorno aos estudos ; estímulo à participação das empresas, por meio da solicitação de recursos materiais, espaço físico para as salas, auxílio financeiro ; incremento dos cursos profissionalizantes no país ; e ampliação da oferta de ensino a distância em todo o território nacional.
Muitos estudantes abordaram, com propriedade, a associação entre um mercado de trabalho cada vez mais competitivo e a necessidade de qualificação profissional, inclusive com uma clara percepção de que os salários são tanto maiores quanto maior a quantidade de anos de estudo e/ou a qualidade da formação.
Nesse tópico, retorna -se ao papel crucial da EJA como viabilizadora dessa formação.
Em continuidade à análise, serão apresentados aspectos negativos que foram notados com frequência nas respostas à questão 10.
De um ponto de vista mais global, a dificuldade repousa no entendimento da questão, isto é, na decifração do comando.
Muitos fugiram completamente do assunto proposto, ainda que o enunciado da questão fosse farto em dados específicos ao tema.
Em vez de responderem diretamente quais eram os problemas enfrentados por aqueles que têm de trabalhar e estudar ao mesmo tempo, muitos estudantes optavam por falar das dificuldades que existem na educação brasileira, de forma estendida e pasteurizada.
Vários estudantes não conseguiam, em suas respostas, relacionar as dificuldades de conciliação da escola com o trabalho, a não ser, como já mencionado, nos casos de depoimentos pessoais, relatos de experiência, ou seja, traços de uma biografia.
Assim, eram respostas cujo teor era apenas o de registrar o abandono dos estudos por parte de vários jovens que têm de trabalhar.
O tom das respostas, nesses casos, era meramente de uma descrição e não de uma análise, muito menos de um elemento que compusesse uma argumentação.
A dificuldade de compreensão da questão se expressava nas respostas que, em lugar de mostrarem os problemas de conciliação entre escola e trabalho, prendiam -se à ideia de “ associação ” – termo que consta do enunciado da questão – entre escola e trabalho.
Muitos alunos relatavam as consequências positivas de uma boa educação no momento em que se busca bons postos de trabalho.
Tratava -se de uma fuga parcial ao comando da questão, evidenciando má compreensão do mesmo.
Uma dificuldade facilmente percebida nas respostas refere -se à tendência de os estudantes tratarem dos problemas relativos à escola e, muito separadamente, daqueles problemas relativos ao trabalho, reconhecendo uma espécie de sinergia que impediria de visualizá -los em sua individualidade, ao mesmo tempo que em sua complementaridade.
Isso significa que os estudantes distinguiam por mero isolamento, ou seja, eles não conseguiam distinguir e articular e, sobretudo, reconhecer as convocações mútuas entre os problemas da escola e do trabalho, no caso explícito dos estudantes trabalhadores.
Essa dificuldade de identificar, associar e analisar aspectos distintos, até mesmo concorrentes, antagônicos e complementares, denuncia a fragilidade de um pensamento simplificador que a qualidade das respostas revelou ser ainda hegemônico nas universidades brasileiras.
Notou -se a presença de preconceitos na análise dos fatores que levam ao abandono dos estudos.
Muitas respostas apontavam como causa desse abandono e da repetência a preguiça, a falta de perseverança e o cansaço dos estudantes.
Nesse tópico específico, fica evidente a responsabilização do indivíduo pelos percalços e fracassos de sua vida, descontextualizando -o de uma sociedade que, em seu movimento de conjunto, lhe impõe condicionantes.
A precariedade do uso da língua pátria de modo culto foi uma constante, conforme mencionado na análise da questão 9.
As impropriedades lexicais e a falta de elementos coesores que articulassem as partes constitutivas das respostas foram um incômodo constante.
Nesse sentido, agrava -se a situação da baixa qualidade textual em decorrência complementar da construção de textos meramente descritivos e não argumentativos, como é exigido no suporte da questão.
Isso demonstra a dificuldade dos estudantes universitários na compreensão de que uma argumentação requer, imperiosamente, justificações para os apontamentos apresentados.
Os textos das respostas eram, em sua maioria, opinativos, sem qualquer compromisso com elementos explicativos ou com a busca de justificativas. 3.3.2 Componente de Conhecimento Específico Na parte da prova relativa ao componente de Conhecimento Específico, observa -se que a diferença de desempenho foi mais acentuada ao se comparar Ingressantes e Concluintes.
A Tabela 3.15 contém dados elucidativos, na qual é possível observar que, no grupo dos Ingressantes, a nota média 4,3 é significativamente inferior à alcançada por Concluintes, 25,6.
Essas médias foram menores do que as obtidas no conjunto das questões objetivas, nas quais Ingressantes alcançaram nota média 34,9, e Concluintes, 52,6 ( Tabela 3.8 ).
O Gráfico 3.12 apresenta a distribuição de notas nos diferentes intervalos.
Observa -se que a faixa modal foi a de notas entre 0 e 10 pontos, tanto para Ingressantes ( 83,9 \% ) quanto para Concluintes ( 27,5 \% ).
Para Ingressantes as proporções decaem para as faixas seguintes, e para Concluintes nota -se uma segunda moda na faixa de 20 a 30 pontos com 18,2 \% dos estudantes deste grupo, mas também uma tendência de queda.
A análise de cada uma destas questões será feita a seguir. ∗ ∗ A análise completa, com apresentação de trechos de respostas de estudantes, pode ser encontrada no documento “ Relatório da Correção das Questões Discursivas ”, elaborado pelos docentes que participaram do processo de correção das mesmas. 3.3.2.1 Análise da Questão Discursiva 38 do Componente de Conhecimento Específico Na questão 38, cujos resultados aferidos encontram -se descritos na Tabela 3.16, o grupo dos Ingressantes obteve nota média 0,9, enquanto a média dos Concluintes foi 17,3.
Deixaram a questão em branco 67,5 \% dos Ingressantes e 19,3 \% dos Concluintes.
Ressalta -se também a elevada proporção de notas zero para aqueles que não deixaram a questão em branco : 26,1 \% dos Ingressantes e 25,1 \% dos Concluintes.
Tanto para Ingressantes quanto para Concluintes as proporções decaem para as faixas seguintes.
A nota máxima, 100,0 pontos, foi alcançada apenas no grupo de Concluintes. 3.3.2.2 Comentários sobre as respostas à Questão Discursiva 38 A questão compreende três subitens que exploram a capacidade de aplicação de conhecimentos adquiridos em disciplinas básicas - Bioquímica e Fisiologia - e conhecimentos específicos abordados em disciplinas profissionalizantes - Patologia Clínica e Clínica Médica.
Em relação às respostas obtidas no subitem “ a ”, percebeu -se a dificuldade de compreensão daquilo que deveria ser respondido, pois, quando perguntado sobre “ os achados esperados no eritrograma ”, os estudantes entenderam que os parâmetros que compõem o eritrograma ( hemoglobinometria, volume globular ou hematócrito, hematimetria, por exemplo ) deveriam ser relatados na questão.
Logo, em várias respostas, foi comum a presença dos parâmetros em substituição aos termos que deveriam ser utilizados ( anemia normocítica e normocrômica, hemoglobina e VGM com valores abaixo do padrão normal ).
Por exemplo, utilizou -se o termo “ Volume Globular Baixo ”, para fazer referência à anemia esperada no resultado do exame.
Em algumas respostas, os estudantes citaram a desidratação como causa de um quadro de policitemia, demonstrando uma inabilidade em realizar a interpretação do caso clínico e aplicando uma espécie de “ memorização ” da causa mais comum de policitemia encontrada em veterinária, sem associar com o contexto da situação, ou seja, não conseguiram desenvolver suas competências e habilidades na solução de casos propostos.
Por outro lado, foi possível observar respostas que demonstraram uma boa capacidade de interpretação.
Muitos associaram a insuficiência renal à falta de produção de eritropoietina, determinando o aparecimento de uma anemia normocítica normocrômica arregenerativa, sendo este o padrão de resposta esperado para tal item.
No subitem “ b ”, a pergunta possui clareza e objetividade.
Entretanto, o despreparo dos estudantes em descrever sobre as alterações esperadas no exame urinálise, evidenciou a ausência de conhecimentos sobre quais são os parâmetros que compõem o exame de urina.
A ausência de capacidade de raciocínio em situações simples, como sobre a função renal, ficou fortemente marcada.
Na maioria das provas corrigidas, a resposta foi “ elevação dos níveis de uréia e creatinina ”, sugerindo que esses estudantes receberam conhecimento e memorizaram que a ureia e a creatinina são parâmetros que avaliam a função renal.
No entanto, como ureia e creatinina são produtos que devem ser excretados pelos rins, a elevação dessas substâncias no sangue caracteriza uma falha no funcionamento renal e não a elevação dessas mesmas substâncias na urina, o que demonstra dificuldade nos conhecimentos em patologia clínica ( nível profissionalizante ) e, principalmente, nos conhecimentos em fisiologia ( nível básico ).
O subitem “ b ” possui grau de dificuldade relativamente baixo, porém, contraditoriamente, foi observado como sendo o de menor índice de acerto.
Percebeu -se, por outro lado, que há memorização de conceitos, mas uma certa dificuldade quanto à aplicação do conhecimento.
No subitem “ c ”, alguns estudantes usaram os conhecimentos de forma adequada, unindo as informações obtidas no ciclo básico àquelas obtidas no ciclo profissional.
Muitos associaram a insuficiência renal com a hipoproteinemia, que é responsável pela diminuição da pressão oncótica e, consequentemente, pelo aparecimento dos edemas.
Em contrapartida, foi possível notar a dificuldade de interpretação de outro grupo de estudantes, que pensou em hipóteses absurdas para justificar a situação apresentada.
Em nenhum momento, o problema expõe o aparecimento de uma hepatopatia, ou de uma insuficiência cardíaca no paciente estudado, mas alguns estudantes justificaram a ascite e o edema do animal com base na “ possibilidade ” de se encontrarem essas alterações, parecendo ser mais fácil trabalharem com hipóteses do que com situações concretas, uma vez que o diagnóstico já havia sido informado.
O número de estudantes com nível de conhecimento intermediário foi muito pequeno.
Os erros de grafia e de concordância estavam presentes de maneira massiva.
A quantidade elevada de questões em branco pode sugerir que os estudantes não cursaram as disciplinas específicas e devem ter optado por não escrever em função de não terem elementos suficientes. 3.3.2.3 Análise da Questão Discursiva 39 do Componente de Conhecimento Específico A Tabela 3.17 contém as informações relativas à questão 39 do conjunto de questões do componente de Conhecimento Específico.
O desempenho dos dois grupos – Ingressantes e Concluintes – nessa questão, foi pouco superior ao desempenho na questão de número 38.
As médias foram 6,9 no grupo dos Ingressantes e 25,2 no dos Concluintes.
A proporção de estudantes que deixaram a questão em branco também foi elevada, sendo os percentuais de 68,7 \% entre os Ingressantes e de 29,4 \% entre os Concluintes.
A nota máxima ( 100,0 ) foi atingida pelos dois grupos.
A proporção de notas zero mesmo entre os que não deixaram a questão em branco foi significativa : 11,7 \% para ambos os grupos.
As faixas de 10 a 20 e de 30 a 40 pontos correspondem às notas de maior frequência nos dois grupos : respectivamente 4,6 e 3,4 \% para Ingressantes e 10,5 e 9,9 \% para Concluintes.
As distribuições das notas entre os estudantes constam do Gráfico 3.14. 3.3.2.4 Comentários sobre as respostas à Questão Discursiva 39 A questão compreende dois subitens que exploram a capacidade de aplicação de conhecimentos abordados em disciplinas profissionalizantes relacionadas com produção animal, zootecnia, medicina veterinária preventiva e saúde pública.
No subitem “ a ”, verificou -se a dificuldade de o estudante de Medicina Veterinária fazer cálculos.
Alguns fizeram o cálculo incompleto - 100 matrizes X 150 litros por matriz - pois não levaram em conta os dias do ano.
Outros multiplicaram 100 matrizes X 150 litros por matriz X 360 dias do ano, em vez de 365 dias, ou seja, não sabiam claramente quantos dias formam um ano.
Um terceiro grupo simplesmente não montou a equação, obtendo resultados bem diferentes.
Houve, ainda, casos em que o raciocínio estava correto, mas o cálculo estava equivocado.
O subitem continha solicitação relativamente simples, tendo ficado evidente a incapacidade de vários dos estudantes em montar a pequena equação solicitada.
No subitem “ b ”, a totalidade das respostas corretas e das parcialmente corretas, os estudantes sugeriram o uso do biodigestor, mas apenas uma pequena parte complementou que, além da produção de energia, este também evita a presença de dejetos no meio ambiente.
A utilização de esterqueira foi uma resposta bastante frequente, poucos, porém, indicaram que a mesma deveria ser vedada para evitar contaminação do solo e águas.
Algumas respostas continham a avaliação do solo e dos dejetos.
A questão foi formulada com clareza e continha informações pertinentes para a elaboração da resposta, alguns estudantes, contudo, \textbf{julgaram} faltar dados para resposta, demonstrando despreparo.
Outros escreveram “ ainda não aprendi na faculdade ”.
Esses, provavelmente, constituem o grupo de iniciantes.
Mesmo não tendo aprendido o conteúdo seria possível, no entanto, responder a parte da questão, especialmente ao subitem “ a ”, que dependia de interpretação e raciocínio lógico-matemático.
Os erros ortográficos estavam presentes em grande quantidade, como na grafia das palavras “ utilisação ”, “ seguranssa ”, dentre outras.
Tais dificuldades, lamentavelmente, tornaram as respostas, em várias circunstâncias, incompreensíveis.
A quantidade de questões em branco pode estar relacionada com o fato de os estudantes não terem cursado todas as disciplinas específicas, embora, diante do universo de respostas apresentadas, tenha ficado evidente que os estudantes de Medicina Veterinária não estão preparados para questões de conhecimento em produção animal, medicina veterinária preventiva e saúde pública. 3.3.2.5 Análise da Questão Discursiva 40 do Componente de Conhecimento Específico A Tabela 3.18 contém as informações relativas à questão 40 do conjunto de questões do componente de Conhecimento Específico.
O desempenho dos Concluintes, nessa questão, foi superior ao obtido nas duas outras questões discursivas desse componente : obtiveram média 34,3, comparando -se aos 17,3 da questão 38 e 25,2 da 39.
A nota média de Ingressantes foi 5,0.
A mediana total foi zero.
Por outro lado, a nota máxima ( 100,0 ) foi obtida nos dois grupos, com os seguintes percentuais : 0,3 \% ( Ingressantes ) e 7,2 \% ( Concluintes ).
A questão foi deixada em branco por 67,3 \% dos Ingressantes e por 15,2 \% dos Concluintes.
A nota zero foi alocada a 20,8 \% dos Ingressantes e a 14,7 \% dos Concluintes.
Entre concluintes, o percentual de 15,9 \% - o mais elevado - se repetiu nas faixas de notas entre 10 e 20 e entre 30 e 40 pontos.
A distribuição de notas dos dois grupos pode ser observada no Gráfico 3.15. 3.3.2.6 Comentários sobre as respostas à Questão Discursiva 40 Embora o enunciado da questão solicitasse que um texto fosse redigido, quatro subitens relativos aos conhecimentos de inspeção de produtos de origem animal, de doenças infecciosas e saúde pública foram determinados, sugerindo que a resposta fosse objetiva.
Diante disso, poucos estudantes desenvolveram um texto.
O subitem “ a ” teve número expressivo de acerto.
Os estudantes que não acertaram o diagnóstico consideraram como resposta doenças que não correspondem aos dados fornecidos no enunciado, tais como : linfadenite caseosa, cisticercose, febre aftosa, cirrose, doença da vaca louca, pneumonia, brucelose, carbúnculo, carcinoma, verminose, inflamação, anemia infecciosa equina, dermatite alérgica a pulga.
Alguns estudantes não conseguiram acertar o diagnóstico, mencionando que o animal deveria ter ingerido ou inalado areia.
Outros escreveram que o animal havia recebido muita pancada antes do abate, justificando a presença dos nódulos.
O subitem “ b ” pode ser considerado como o de maior grau de dificuldade.
Embora preciso na solicitação, deixou o estudante em dúvida, pois a maioria acabou escrevendo sobre a fisiopatogenia da tuberculose, ou sobre todas as formas de infecção para o homem e para os animais.
Poucos associaram as lesões restritas ao sistema respiratório e linfonodos regionais como indicativo do modo de transmissão por inalação do agente.
Outro erro que evidencia a má interpretação do texto foram as respostas sobre a doença na propriedade, quando a informação versava sobre a situação em um matadouro.
Os subitens “ c ” e “ d ” foram os que obtiveram maior quantidade de respostas erradas, apesar de não apresentarem solicitações difíceis.
Pouquíssimos estudantes destinaram as carcaças caquéticas e os órgãos que entraram em contato com o líquido caseoso à condenação total e à condenação, respectivamente.
Muitos estudantes acusaram o descarte das estruturas para graxaria, ou para incineração, mas sem definição de condenação.
Alguns sugeriram um “ banho a vapor ” para esterilização.
Houve quem destinasse as carcaças para produção de embutidos, de carne seca e de farinha.
Respostas bem escritas foram muito poucas nesta questão.
Contava, principalmente, com erros gramaticais, palavras com a escrita incorreta ( gracharia, icineração, ensineração, exumeração ) ausência de concordância verbal e textos incompreensíveis.
Poucos estudantes redigiram respostas perfeitas, inclusive reconhecendo que animais caquéticos têm abate diferenciado.
A quantidade de questões em branco pode estar relacionada com o fato de os estudantes não terem cursado as disciplinas específicas.
CONSIDERAÇÕES FINAIS A questão 38, principalmente, possibilitou a avaliação dos conhecimentos do ciclo básico associados àqueles do ciclo profissionalizante, e se percebeu que o nível de conhecimento dentro do assunto abordado está muito aquém do esperado.
Falta conhecimento básico sobre termos técnicos utilizados na patologia clínica.
Foi muito frequente o aluno não ter noção dos achados do eritrograma, da urinálise e da bioquímica sérica.
Os estudantes demonstraram que os conhecimentos do ciclo básico do curso tiveram uma assimilação bem melhor do que os conhecimentos do ciclo profissional.
Isso ficou claro pelo número de estudantes que acertou o item ( c ) da questão ( onde o conhecimento de osmolaridade sanguínea auxiliava a resposta, conteúdo que foi visto em fisiologia e/ou em bioquímica, por exemplo ), os mesmos estudantes, porém, não tinham noção do que era perguntado nos itens “ a ” e “ b ”.
Na questão 39, houve grande quantidade de respostas totalmente sem propósito, citam -se, dentre elas, o desprezo dos dejetos diretamente nos rios, sem tratamento prévio, o incentivo à diminuição do consumo de carne suína, o incentivo ao ” vegetarianismo ”, a diminuição do número de matrizes ( já que são elas as ” grandes causadoras ” do problema ) e diminuição do número de animais por granja, mostrando que muitos estudantes desconhecem assuntos de alta relevância ( sustentabilidade ou preservação de recursos, por exemplo ) e que são discutidos no mundo todo como fundamentais para a saúde do planeta.
A questão 40 pôde demonstrar que os estudantes ignoram a importância da tuberculose e as determinações do RIISPOA ( Regulamento de Inspeção Industrial e Sanitária de Produtos de Origem Animal ) para essa condição, deflagrando o despreparo de futuros profissionais, que reconhecem e liberam as carcaças portando lesões, reconhecidamente tuberculosas ou não, para consumo humano e de outros animais.
Em alguns casos, percebe -se que o estudante não conhece o significado das palavras ( “ doença patológica ” ; “ O parasita da lesão em questão se apresenta muito volátil ” ; “... a cremação com posterior exumeração dos restos ” ).
Outros estudantes assumiram desconhecer os assuntos por não se interessarem pelas áreas em que os mesmos se encaixam.
A qualidade das respostas ofereceu à banca de corretores uma ampla visão acerca da graduação em Medicina Veterinária, habilitando -a a concluir que há necessidade premente de promover melhorias não apenas na educação superior da área, mas fundamentalmente na educação básica, que fornece os alicerces para o posterior desenvolvimento de estudantes/profissionais com competências e habilidades bem estruturadas, e que, por isso, poderão desempenhar suas funções com êxito.

% matched lemmas: julgar
\end{textsample}
